\section{Worked examples}\label{sec:worked}

The examples below take one step of the method at a time and carry it through
with the numbers of this paper, so that an analyst can reproduce each figure by
hand. Every number comes from the runs reported in Sections \ref{sec:eval} and
\ref{sec:security}, and \nolinkurl{scripts/worked_examples.py} recomputes all of
them into \nolinkurl{artifacts-robustness/worked_examples.json}.
Examples \ref{ex:bound} and \ref{ex:selection} show how a floor is decided,
Example \ref{ex:history} shows the attack that the floor must resist, Examples
\ref{ex:zero} to \ref{ex:length} show how many attacks an audit costs, and
Examples \ref{ex:grinding} and \ref{ex:dp} show what can go wrong and how the
result compares with differential privacy.

\begin{example}[One coordinate, one bound]\label{ex:bound}
Take the single-release attack with one exposure reading at $k_{\min} = 1$, the
first row of Table \ref{tab:robust}. On the confirmation draw the attack named
the right loan for 337 of 4{,}000 targets, a rate of 0.084. Rule \eqref{eq:admiss}
compares the exact one-sided 95 per cent upper bound with the tolerance; that
bound is the 0.95 quantile of the $\mathrm{Beta}(337+1,\,4000-337)$
distribution, which is 0.092. The tolerance is
0.12, so the bound meets it with a margin of $1 - 0.092/0.12 = 0.23$, and the
file passes for this adversary with no suppression. The decision uses the bound,
not the rate: had the rate been 0.115 with the same number of targets, the bound
would have been 0.123 and the release would have failed, as the next example shows.
\end{example}

\begin{example}[Choosing a floor when the confirmation fails]\label{ex:selection}
For the adversary who pools six vintages, holds one reading and knows the
target's history, Algorithm \ref{alg:selection} tests the floors from 300 down on
the selection draw of 4{,}000 targets each. The upper bounds are 0.042 at 300,
0.059 at 200, 0.073 at 150, 0.081 at 100, 0.101 at 75 and 0.111 at 50, all below
0.12, and 0.139 at 40, which fails. The search stops there and selects 50, the
last floor that passed. On the independent confirmation draw the same attack at
50 succeeded for 458 of 4{,}000 targets, a rate of 0.115 with a bound of 0.123,
above the tolerance. The algorithm therefore moves to the next stricter floor,
75, where 384 of 4{,}000 targets were named, a rate of 0.096 with a bound of
0.104, which passes. The floor for this adversary is 75, as Table
\ref{tab:robust} reports. Falling back to 75 keeps the error bound, because 75
had already passed at selection; choosing 40 or 30 after seeing these numbers
would not.
\end{example}

\begin{example}[The history attack on one cell]\label{ex:history}
Consider the four loans of Figure \ref{fig:history}, all in the target's key cell,
with published log exposures 12.0 for the target $t$, 12.7 for $a$, 12.25 for
$b$ and 11.6 for $c$. The adversary's reading of the target is 12.30, with
standard deviation 0.6. Without the history, Algorithm \ref{alg:attack} scores
each loan by the normal density of the reading, and the normalised posteriors are
0.277 for $t$, 0.251 for $a$, 0.313 for $b$ and 0.159 for $c$. The adversary names
$b$, which is wrong, and with a posterior below one half it would not act on the
claim. With the history, $b$ and $c$ are excluded, because $b$ is live in a
vintage where the target is not and $c$ is absent from one where the target is
live. Only $t$ and $a$ remain, with posteriors 0.524 and 0.476. The adversary now
names the target and, with a posterior above one half, acts. Nothing in any
single release changed; the gain comes entirely from the pattern of appearances
across vintages, which is why pooled attacks need coordinates of their own.
\end{example}

\begin{example}[How many attacks certify a coordinate with no successes]\label{ex:zero}
Proposition \ref{prop:cert} gives the fixed-sample count: with no success in $B$
attacks the bound meets $\tau$ when $B \ge \log\alpha/\log(1-\tau)$. At
$\tau = 0.12$ and $\alpha = 0.05$ this is $2.996/0.128 = 23.4$, so 24 attacks. If
the coordinate carries only a third of the error budget, as each attack does when
three are split equally, $\alpha$ becomes $0.05/3$ and the count rises to
$\lceil 4.094/0.128\rceil = 33$. Splitting the budget is therefore cheap when the
risk is far below the tolerance: the count grows with the logarithm of the number
of coordinates, not in proportion to it. The release chain of Algorithm
\ref{alg:chain} avoids even this cost, because a floor that fails ends the search
and each attack keeps the full level.
\end{example}

\begin{example}[The betting gate, draw by draw]\label{ex:gate}
In Algorithm \ref{alg:chain} a coordinate certifies when its evidence $E$ reaches
$1/\alpha = 20$. Fix the bet at $q = 0.05$ with $\tau = 0.12$. Each failed attack
multiplies $E$ by $(1-q)/(1-\tau) = 0.95/0.88 = 1.080$, and each success
multiplies it by $q/\tau = 0.05/0.12 = 0.417$. Starting from $E = 1$, a run of
failures reaches 20 after $\log 20/\log 1.080 = 39.1$ draws, so on the 40th. One
success sets the evidence back by a factor of 0.417, which takes
$\log(1/0.417)/\log 1.080 = 11.4$ further failures to recover. An attack that
succeeds on one target in twenty gains about 0.029 in $\log E$ per draw and so
certifies after about a hundred draws on average, while one that succeeds on one
in nine never gains ground: at
$p = 0.116$ the expected change in $\log E$ per draw is negative for this bet.
The auditor may look at $E$ after every draw and stop at the first crossing
without any correction, which a fixed-sample bound does not allow.
\end{example}

\begin{example}[Why the release chain stops at 150]\label{ex:length}
Near the tolerance the cost of certification is set by the gap between the true
risk and $\tau$. With the best fixed bet, the expected number of draws is about
$\log(1/\alpha)$ divided by the Kullback--Leibler divergence between
$\mathrm{Bernoulli}(p)$ and $\mathrm{Bernoulli}(\tau)$ (the expected-length
proposition of the companion article). For the full pooled attack at
$k_{\min} = 150$, where the true risk is 0.091, that is about 720 draws, well
within the cap of 3{,}000. At $k_{\min} = 100$ the true risk is 0.116, only 0.004
below the tolerance, and the same calculation gives about 43{,}000 draws. The
release chain of Figure \ref{fig:chain} therefore certifies 150 in almost every
stream and 100 only by chance. The floor of 100 is safe, but an audit of any
practical size cannot show it.
\end{example}

\begin{example}[What seed grinding buys an operator]\label{ex:grinding}
Table \ref{tab:audit-process} reports that an honest audit of a coordinate whose
true risk equals the tolerance certifies it falsely in 3.8 per cent of runs. An
operator who can try ten seeds and report the first that passes succeeds with
probability $1 - (1 - 0.038)^{10} = 0.32$, and with a hundred seeds
$1 - (1-0.038)^{100} = 0.98$. The stated error level of 0.05 then means nothing.
Algorithm \ref{alg:committed} removes the choice: the targets are fixed by a
secret committed at registration and by a beacon value published later, so the
operator has one stream per coordinate and the rate returns to 0.038.
\end{example}

\begin{example}[The same tolerance under differential privacy]\label{ex:dp}
Suppose instead that a release were $(\varepsilon,\delta)$-differentially private
with $\varepsilon = 1$ and $\delta = 10^{-6}$, and that a declared attack succeeds
3 per cent of the time when the target's record is replaced by another. The
hypothesis-testing bound of Section \ref{sec:dp} caps that attack's success at
$e^{1}\times 0.03 + 10^{-6} = 0.082$, below the tolerance of 0.12, and the same
cap holds for every attack at once, declared or not. The gate cannot give that
guarantee: its certificate covers the attacks that were registered. What it can
do is certify releases, such as the suppressed loan-level files of this paper,
that have no privacy parameter at all.
\end{example}
